% from ega0-19b-fr.tex
\oldpage[0\textsubscript{IV-1}]{114}

par $u$ de l'élément $p\cdot1$ de $W$ est l'élément $p\cdot1$ de $A$~; comme
l'élément $p\cdot1$ de $W$ engendre l'idéal maximal de cet anneau, on déduit
aussitôt de l'hypothèse que
$\gr(u):\gr_\bullet(W)\to\gr_\bullet(A)$ est surjectif, et par suite il en est
de même de $u$ (Bourbaki, \emph{Alg. comm.}, chap.~III, §~2,
n\textsuperscript{o}~8, cor.~2 du th.~1).
\end{proof}

\subsection{Algèbres relativement formellement lisses}
\label{subsection:0.19.9-fr}

\begin{definition}[19.9.1]
\label{0.19.9.1-fr}
Soient $\Lambda$ un anneau topologique, $A$ une $\Lambda$-algèbre topologique,
$B$ une $A$-algèbre topologique. On dit que $B$ est une $A$-algèbre
\emph{formellement lisse relativement à $\Lambda$} si, pour toute $A$-algèbre
topologique discrète $C$, et tout idéal nilpotent $\mathfrak J$ de $C$, tout
$A$-homomorphisme continu $u_0:B\to C/\mathfrak J$ qui se factorise en
$B\xrightarrow{u}C\to C/\mathfrak J$, où $u$ est un $\Lambda$-homomorphisme
continu, se factorise aussi en $B\xrightarrow{v}C\to C/\mathfrak J$, où $v$
est un $A$-homomorphisme continu.
\end{definition}

Il résulte de cette définition que si $B$ est une $A$-algèbre \emph{formellement
lisse}, alors $B$ est aussi \emph{formellement lisse relativement à
$\Lambda$}, pour toute structure de $\Lambda$-algèbre topologique définie sur
$A$ (autrement dit, tout homomorphisme continu d'anneaux $\Lambda\to A$).

\begin{proposition}[19.9.2]
\label{0.19.9.2-fr}
Soient $\Lambda$ un anneau topologique, $A$ une $\Lambda$-algèbre topologique.
\begin{enumerate}
  \item[(i)] $A$ est une $\Lambda$-algèbre formellement lisse relativement à
  $\Lambda$.
  \item[(ii)] Si $B$ est une $A$-algèbre formellement lisse relativement à
  $\Lambda$ et $C$ une $B$-algèbre formellement lisse relativement à
  $\Lambda$, alors $C$ est une $A$-algèbre formellement lisse relativement à
  $\Lambda$.
  \item[(iii)] Soient $B$ une $A$-algèbre formellement lisse relativement à
  $\Lambda$, $A'$ une $A$-algèbre topologique, alors la $A'$-algèbre
  topologique $B\otimes_A A'$ est formellement lisse relativement à $\Lambda$.
  \item[(iv)] Soient $B$ une $A$-algèbre topologique, $S$ (resp. $T$) une
  partie multiplicative de $A$ (resp. $B$) telle que l'image canonique de $S$
  dans $B$ soit contenue dans $T$. Si $B$ est une $A$-algèbre formellement
  lisse relativement à $\Lambda$, alors $T^{-1}B$ est une $S^{-1}A$-algèbre
  formellement lisse relativement à $\Lambda$.
  \item[(v)] Soient $B_i$ ($1\leq i\leq n$) des $A$-algèbres topologiques.
  Pour que $\prod_{i=1}^n B_i$ soit une $A$-algèbre formellement lisse
  relativement à $\Lambda$, il faut et il suffit que chacune des $B_i$ le
  soit.
\end{enumerate}
\end{proposition}

\begin{proof}
L'assertion (i) est triviale, et la démonstration des autres est étroitement
calquée sur les démonstrations de (19.3.5)~; elle est donc laissée au lecteur.
\end{proof}

\begin{corollary}[19.9.3]
\label{0.19.9.3-fr}
Soient $A$ un anneau topologique, $A'$ et $B$ deux $A$-algèbres topologiques.
Alors la $A$-algèbre topologique $B'=B\otimes_A A'$ est une $B$-algèbre
formellement lisse relativement à $A$.
\end{corollary}

\begin{proof}
Cela résulte de (19.9.2, (i) et (iii)).
\end{proof}

\begin{proposition}[19.9.4]
\label{0.19.9.4-fr}
Soient $\Lambda$ un anneau topologique, $A$ une $\Lambda$-algèbre topologique,
$B$ une $A$-algèbre topologique. Les conditions suivantes sont équivalentes~:
\begin{enumerate}
  \item[a)] $B$ est une $A$-algèbre formellement lisse relativement à
  $\Lambda$.
  \item[b)] $\widehat B$ est une $A$-algèbre formellement lisse relativement à
  $\Lambda$.
  \item[c)] $\widehat B$ est une $\widehat A$-algèbre formellement lisse
  relativement à $\Lambda$.
  \item[d)] $\widehat B$ est une $\widehat A$-algèbre formellement lisse
  relativement à $\widehat\Lambda$.
\end{enumerate}
\end{proposition}

\begin{proof}
On laisse encore au lecteur la démonstration, calquée sur celle de (19.3.6).
\end{proof}

